\begin{exo}
	Le plan est muni d'un repère orthonormé $\OIJ$ d'unité $2\text{ cm}$.

	On considère les points $A\coordl{2}{1}$, $B\coordl{5}{1}$, $C\coordl{5}{-2}$ et $D\coordl{2}{-2}$.
	\begin{enumerate}
		\item Faire une figure.
		\item
		      \begin{enumerate}
			      \item Déterminer les coordonnées de $K$, milieu de $[AC]$.
			      \item Déterminer les coordonnées de $L$, milieu de $[BD]$.
			      \item En déduire que $ABCD$ est un parallélogramme.
		      \end{enumerate}
		\item
		      \begin{enumerate}
			      \item Calculer les distances $AC$, $AD$ et $DC$.
			      \item En déduire la nature du triangle $ADC$.
		      \end{enumerate}
		\item Conclure sur la nature du parallélogramme $ABCD$.
	\end{enumerate}
\end{exo}

